\section{Open problems}\label{op:section}
The equality $\operatorname{gap}(K)=2/n$ is proved here for $d\in\{1,n-2\}$
and in the cases of Theorem~\ref{thm:small}. It is not proved for any other
regular degree sequence. For $n=12$ and $d\in\{5,6\}$ there are $7{,}849$
classes and $83{,}858$ rooted classes; a floating-point computation gives
$\mu\approx14.62>11$, which is evidence and not a proof.

\begin{conjecture}[Regular Curveball gap]\label{op:conjecture}
For every $n\ge4$ and $1\le d\le n-2$ with $nd$ even,
$\operatorname{gap}(H)=n-1$ and $\operatorname{gap}(K)=2/n$ on $\Omega_{n,d}$.
\end{conjecture}

By Theorem~\ref{thm:main}(3), Conjecture~\ref{op:conjecture} is equivalent to
$\mu\ge n-1$, an inequality for functions of the rooted isomorphism class. On
the centered functions of the unrooted isomorphism class, which form
$U_{(n)}\cap\mathbf1^\perp$, the comparison of Proposition~\ref{cmp:main}
gives no information, and the only bound proved for general $n$ and $d$ is
$H\succeq2I$~\cite[Theorem~1.1]{S26}. On $U_{(n-1,1)}$ the proved bound is
$n/2$. A lower bound for $\mu$ of order $n$ would already show that the
relaxation time of $K$ on regular graphs has order $n$. By
\eqref{exp:target} the conjecture is also equivalent to
$S+D\succeq(n-2)H$, and Proposition~\ref{exp:obstruction} shows that this
cannot be obtained from separate bounds on $S$ and $D$. The conjecture has the
form of the statement proved by Caputo, Liggett and
Richthammer~\cite{CLR10} for the interchange process: the gap of the process
is attained on a small explicit invariant subspace, here the span of the
edge indicators. We do not know whether their method applies; conditioning
on the neighborhood of one vertex leaves the class of regular graphs. A closer
analog is the Kac walk and its relatives, where the gap is attained at an
explicit function and is obtained from the three-particle gap by the
expansion of Section~\ref{exp:section}~\cite{Caputo08}; for regular graphs
Proposition~\ref{exp:obstruction}(1) shows that the three-vertex operators do
not have the gap this route requires.

In Table~\ref{cert:table} the computed values of $\mu$ exceed $n-1$ by at
least $1.2$, and the least ratio $\mu/(n-1)$ is about $1.17$, at $n=8$,
$d=3$. We do not know whether $\mu-(n-1)$ is bounded below by a positive
constant, or grows with $n$, over all regular degree sequences.

For a degree sequence with unequal degrees, Remark~\ref{cmp:general} compares
only the trades at pairs of equal degree with transpositions. No analog of
Theorem~\ref{thm:main}(1) is known there, and the gap of $H$ can equal
$2$~\cite[Theorem~1.1]{S26}. For binary matrices with equal row sums,
Proposition~\ref{bin:main} leaves the row-symmetric functions. The
margins for which $\mu_{\mathrm{sym}}\ge m/2$ are not characterized; the
criterion \eqref{bin:caputo} is sufficient and is not necessary.

The bound \eqref{eq:mixupper} loses the factor $\log M$ of the passage from
the spectral gap to total variation. For $d=1$ the walk is a quotient of the
random-transposition walk on $S_n$ and mixes in order $n\log n$
steps~\cite{DS81,DH02}, while \eqref{eq:mixupper} has order $n^2\log n$. An
entropy or log-Sobolev inequality for $K$ on regular graphs is not proved
here.
